% Chapter 2: device and precision utilities

\chapter{Device and precision utilities}

\label{Chapter2}

The module \code{fluxlb.core.gpu} holds no physics. It exists so that the rest of the
package asks one place for the device, the accumulation dtype and memory or timing figures,
and behaves the same on a CUDA node, an Apple laptop and a plain CPU. Signatures and
per-function behaviour are in the API reference generated from the docstrings at
\url{https://cerg-flux-lab.github.io/fluxlb/}; this chapter explains the policies behind them.

%----------------------------------------------------------------------------------------

\section{Device selection}

\code{get\_device()} returns the first available of CUDA, MPS and CPU. The environment
variable \code{FLUXLB\_DEVICE} overrides the choice (for example \code{cpu} or
\code{cuda:1}) so a SLURM script or a test can pin the device without editing code. An
override that names an unavailable backend raises \code{RuntimeError} rather than falling
back, so a job never quietly runs on the CPU.

\code{resolve\_device(x)} normalises a \code{torch.device}, a device string, a CUDA index
or \code{None} (meaning \code{get\_device()}) to a \code{torch.device}. Every per-device
helper accepts this loose form. \Cref{fig:device} shows the resolution order.

\begin{figure}[h]
\centering
\begin{adjustbox}{max width=\linewidth}% Device resolution in get_device(): override first, then CUDA > MPS > CPU.
\begin{tikzpicture}[
  >={Stealth[length=2mm]},
  node distance=6mm and 10mm,
  dec/.style={draw, diamond, aspect=2.6, align=center, font=\scriptsize, inner sep=1pt, fill=gray!5},
  dev/.style={draw, rounded corners=2pt, align=center, font=\small\ttfamily, inner sep=5pt, minimum width=18mm},
  err/.style={dev, fill=red!8, draw=red!60!black, font=\small},
  arr/.style={->, line width=0.8pt},
  lab/.style={font=\scriptsize, fill=white, inner sep=1pt},
]
\node[dec] (env) {\code{FLUXLB\_DEVICE}\\ set?};
\node[dec, below=of env] (avail) {that backend\\ available?};
\node[dev, fill=green!8, draw=green!50!black, right=14mm of avail] (over) {torch.device(override)};
\node[err, left=14mm of avail] (raise) {RuntimeError};

\node[dec, below=of avail] (cuda) {CUDA\\ available?};
\node[dev, fill=blue!8, draw=blue!60!black, right=14mm of cuda] (dcuda) {cuda};
\node[dec, below=of cuda] (mps) {MPS\\ available?};
\node[dev, fill=blue!8, draw=blue!60!black, right=14mm of mps] (dmps) {mps};
\node[dev, fill=blue!8, draw=blue!60!black, below=of mps] (dcpu) {cpu};

\draw[arr] (env) -- node[lab, right] {yes} (avail);
\draw[arr] (avail) -- node[lab, above] {yes} (over);
\draw[arr] (avail) -- node[lab, above] {no} (raise);
\draw[arr] (env.west) -- ++(-1.1,0) node[lab, above] {no} |- (cuda.west);
\draw[arr] (cuda) -- node[lab, above] {yes} (dcuda);
\draw[arr] (cuda) -- node[lab, right] {no} (mps);
\draw[arr] (mps) -- node[lab, above] {yes} (dmps);
\draw[arr] (mps) -- node[lab, right] {no} (dcpu);

\node[font=\scriptsize, gray!80!black, align=left, anchor=north west, text width=4.2cm]
  at ($(dcuda.east)+(0.4,0.3)$)
  {The override never falls back silently: a SLURM job asking for CUDA on a node without it fails at start-up, not after an hour on the CPU.};
\end{tikzpicture}
\end{adjustbox}
\caption{Device resolution in \code{get\_device()}. An explicit override is honoured or
rejected; only in its absence does automatic selection run.}
\label{fig:device}
\end{figure}

%----------------------------------------------------------------------------------------

\section{Precision policy}\label{sec:precision}

\code{accumulation\_dtype(device)} returns \code{torch.float64} on CUDA and CPU and
\code{torch.float32} on MPS, with a \code{UserWarning} in the MPS case. The fp32 fallback is
intentional: it is the only way to run on MPS at all, and the tests pin the behaviour.

The mechanism for using it is PyTorch's reduction \code{dtype} argument, which upcasts
inside the sum without materialising an fp64 copy of $f$ and is autograd-safe:
\begin{lstlisting}[style=py]
rho = f.sum(dim=0, dtype=accumulation_dtype(f.device)).to(f.dtype)
\end{lstlisting}

%----------------------------------------------------------------------------------------

\section{Memory}

All figures are in bytes and are zero on backends whose allocator PyTorch does not track
(CPU, MPS).
\begin{itemize}[nosep]
\item Allocator counters, as thin wrappers: \code{memory\_allocated(device)},\newline
      \code{peak\_memory\_allocated(device)} and \code{reset\_peak\_memory(device)}.
\item \code{track\_memory(device)}: a context manager that resets the peak counters on
      entry, synchronises on exit, and yields a \code{MemoryReport} with peak allocated, peak
      reserved and net change. Use it around one time step or one training step to size
      gradient checkpointing against \SI{16}{GB} of VRAM.
\item \code{clear\_gpu\_cache()}: runs \code{gc.collect()} first, then releases cached
      allocator memory on CUDA or MPS. It cannot free tensors that are still referenced.
\item \code{format\_bytes(n)}: human-readable rendering, for example \code{15.7 GiB}.
\end{itemize}

%----------------------------------------------------------------------------------------

\section{Timing}

GPU kernels launch asynchronously, so a wall-clock interval taken without synchronising
measures the launch queue, not the computation. \code{synchronize(device)} blocks until
queued work finishes (a no-op on CPU). \code{Timer(device)} is a context manager that
synchronises on entry and exit and exposes \code{elapsed} in seconds.

%----------------------------------------------------------------------------------------

\section{Logging}

SLURM nodes are standalone, so every job log should say where it ran.
\code{device\_summary(device)} returns one line with hostname, Python and torch versions,
and on CUDA the driver version, GPU name, VRAM, compute capability and multiprocessor
count. Print it at job start.
